% =============================================================
%  Relatorio de Pesquisa de Satisfacao - PetCuida (versao Alpha 3i)
%  Registro formalizado de entrevista semiestruturada
%  Compilar com: pdflatex entrevista_petcuida_alpha3i.tex (2x)
% =============================================================
\documentclass[12pt,a4paper]{article}

\usepackage[utf8]{inputenc}
\usepackage[T1]{fontenc}
\usepackage[brazil]{babel}
\usepackage{lmodern}
\usepackage[margin=3cm]{geometry}
\usepackage{setspace}
\usepackage{booktabs}
\usepackage{longtable}
\usepackage{array}
\usepackage{enumitem}
\usepackage{csquotes}
\usepackage{fancyhdr}
\usepackage{titlesec}
\usepackage[hidelinks]{hyperref}

\onehalfspacing

\pagestyle{fancy}
\fancyhf{}
\fancyhead[L]{\footnotesize Pesquisa de satisfação --- PetCuida \textit{Alpha 3i}}
\fancyhead[R]{\footnotesize \thepage}
\renewcommand{\headrulewidth}{0.4pt}

% Ambiente para os turnos de fala da transcricao
\newenvironment{turnos}
  {\begin{list}{}{%
     \setlength{\leftmargin}{2.2cm}%
     \setlength{\labelwidth}{2.0cm}%
     \setlength{\labelsep}{0.2cm}%
     \setlength{\itemsep}{0.6em}%
     \setlength{\parsep}{0pt}}}
  {\end{list}}
\newcommand{\fala}[1]{\item[\textbf{#1}\hfill]}

\newcommand{\entrevistador}{Gabriel}
\newcommand{\participante}{Dona Aline}

\title{\textbf{Avaliação de Satisfação e Usabilidade Percebida do\\
Aplicativo PetCuida (versão \textit{Alpha 3i})}\\[0.4cm]
\large Registro formalizado de entrevista semiestruturada com usuária potencial}
\author{Entrevistador: \entrevistador{} (aluno responsável pela coleta)\\
Participante: \participante{} (usuária potencial)}
\date{Belo Horizonte, 16 de setembro de 2026}

\begin{document}

\maketitle
\thispagestyle{empty}

\begin{abstract}
\noindent Este documento formaliza, para fins acadêmicos, o registro de uma entrevista
semiestruturada realizada com uma usuária potencial do aplicativo \textbf{PetCuida}, em sua
versão \textit{Alpha 3i}. A sessão teve por objetivo apresentar a proposta funcional do
software e coletar a percepção de satisfação, facilidade de uso e interesse de adoção por
parte da participante. Os dados foram obtidos por meio de demonstração guiada do protótipo
funcional em dispositivo Android, seguida de questionamento aberto. Os resultados indicam
avaliação positiva quanto à facilidade de uso percebida, verbalizada pela participante por
meio de analogia com aplicativos de mensageria de uso cotidiano, e manifestação explícita
de interesse no serviço proposto. Não foram registradas sugestões de melhoria nesta rodada,
o que constitui, em si, uma limitação metodológica discutida ao final.

\vspace{0.4cm}
\noindent\textbf{Palavras-chave:} usabilidade percebida; pesquisa de satisfação; entrevista
entrevista semiestruturada; aplicativos móveis; bem-estar animal; PetCuida.
\end{abstract}

\newpage
\tableofcontents
\newpage

% =============================================================
\section{Identificação da Coleta}

\begin{table}[h!]
\centering
\begin{tabular}{@{}>{\raggedright\arraybackslash}p{5.2cm}>{\raggedright\arraybackslash}p{8.5cm}@{}}
\toprule
\textbf{Item} & \textbf{Descrição} \\
\midrule
Objeto avaliado & Aplicativo PetCuida - versão \textit{Alpha 3i} (protótipo funcional Flutter/Android) \\
Instrumento & Entrevista semiestruturada com demonstração guiada do protótipo \\
Entrevistador & \entrevistador{} (aluno responsável) \\
Participante & \participante{} (identificada de forma não nominal ao longo do texto) \\
Perfil do participante & Usuária potencial, sem formação técnica em computação, tutora de animal de estimação \\
Data da coleta & 16 de setembro de 2026 \\
Duração aproximada & 7 minutos e 32 segundos de registro em áudio \\
Local & Ambiente doméstico da participante (entrevista presencial) \\
Registro & Áudio digital (formato \texttt{opus}) e transcrição com marcação temporal (\texttt{srt}) \\
Modalidade de análise & Análise qualitativa de conteúdo, por categorização temática \\
\bottomrule
\end{tabular}
\caption{Ficha técnica da coleta de dados.}
\end{table}

% =============================================================
\section{Objetivo da Pesquisa}

A presente coleta teve como objetivo geral \textbf{verificar a aceitação e a satisfação
percebida de uma usuária potencial diante da proposta funcional do aplicativo PetCuida em
estágio alfa}. Como objetivos específicos, delimitam-se:

\begin{enumerate}[label=\alph*)]
  \item apresentar à participante o conceito do produto e seu modelo de operação;
  \item demonstrar, de forma guiada, o fluxo de telas implementado na versão \textit{Alpha 3i};
  \item registrar a percepção da participante quanto à facilidade de uso do software;
  \item verificar o interesse declarado na contratação ou utilização do serviço proposto;
  \item coletar eventuais sugestões de melhoria e pontos de atrito percebidos.
\end{enumerate}

Ressalta-se que \textbf{esta rodada não constituiu teste de usabilidade com execução de
tarefas pelo usuário}, mas sim apresentação da proposta e coleta de \textit{feedback}
perceptivo, conforme explicitado pelo entrevistador ao encerramento da sessão.

% =============================================================
\section{Descrição do Objeto Avaliado}

O PetCuida é um aplicativo móvel cuja proposta consiste em intermediar, em uma mesma
plataforma, a demanda e a oferta de serviços de cuidado com animais de estimação. O modelo
foi verbalizado pelo entrevistador por meio da analogia de um \enquote{\textit{Uber} de mão
dupla}: o mesmo usuário pode figurar tanto como solicitante quanto como prestador de
serviços, tais como passeio, hospedagem e transporte do animal, enquanto profissionais
da medicina veterinária e clínicas conveniadas podem prestar atendimento pela plataforma.
A contrapartida pela prestação de serviços é convertida em \textbf{crédito} utilizável pelo
próprio usuário para custear os cuidados de seu animal, constituindo o eixo de acessibilidade
econômica da proposta.

A versão \textit{Alpha 3i} implementa, com dados locais e sem dependência de
\textit{backend}, o fluxo completo do diagrama de telas, conforme apresentado na
Tabela~\ref{tab:modulos}.

\begin{table}[h!]
\centering
\small
\begin{tabular}{@{}>{\raggedright\arraybackslash}p{4.2cm}>{\raggedright\arraybackslash}p{9.5cm}@{}}
\toprule
\textbf{Módulo} & \textbf{Funcionalidades demonstradas na sessão} \\
\midrule
Autenticação & Tela de carregamento; acesso por usuário e senha; acesso e registro por conta Google (simulado nesta versão, com retorno visual de confirmação); criação de conta. \\
\textit{Onboarding} do pet & Seleção de espécie (cão, gato, pássaro ou outro); cadastro de nome; campo opcional de raça; seleção de data de nascimento por calendário; definição de foto por galeria ou captura imediata. \\
Área do pet & Painel de cuidados; situação vacinal; calendário de vacinas; agendamentos; prontuário do animal; edição dos dados cadastrados. \\
Triagem de serviço & Pré-seleção orientativa da necessidade do usuário, com possibilidade de interação conversacional. \\
Rede solidária & Listagem de serviços disponíveis; mapa esquemático de pessoas e prestadores próximos; solicitação de serviço. \\
Perfil e configurações & Dados do tutor; definição de foto; informações do programa; termos de uso; políticas de privacidade. \\
\bottomrule
\end{tabular}
\caption{Módulos do PetCuida \textit{Alpha 3i} percorridos durante a demonstração.}
\label{tab:modulos}
\end{table}

Para a demonstração, foi cadastrado um animal fictício de nome \textit{Tobi}, sem raça
definida (SRD), com data de nascimento em setembro de 2010.

% =============================================================
\section{Procedimentos Metodológicos}

\subsection{Abordagem}

Adotou-se abordagem \textbf{qualitativa}, de caráter \textbf{exploratório}, por meio de
entrevista semiestruturada conduzida individualmente. A opção por esse método decorre do
estágio inicial de maturidade do artefato: em versões alfa, interessa menos a mensuração
estatística de desempenho e mais a captura da compreensão, da expectativa e da reação
espontânea do público-alvo diante da proposta de valor.

\subsection{Instrumento e condução}

A sessão foi organizada em três momentos:

\begin{enumerate}[label=\textbf{\arabic*.}]
  \item \textbf{Contextualização} - retomada de conversa anterior com a participante e
        exposição verbal do conceito do produto, de seu modelo de intermediação e do sistema
        de créditos;
  \item \textbf{Demonstração guiada} - percurso narrado pelas telas do protótipo em
        dispositivo móvel, com verbalização, pelo entrevistador, da função de cada elemento
        de interface e das limitações conhecidas do estágio alfa;
  \item \textbf{Coleta de percepção} - aplicação de perguntas abertas relativas à percepção
        geral, ao interesse de adoção e a eventuais sugestões de melhoria.
\end{enumerate}

\subsection{Roteiro aplicado}

\begin{enumerate}[label=\textbf{Q\arabic*.},leftmargin=1.6cm]
  \item Com relação ao que a senhora viu do aplicativo, qual é a sua percepção? O que a
        senhora acha?
  \item Se por acaso fosse oferecido um serviço desse tipo, seria interessante para a senhora?
  \item Havendo algum ponto de vista sobre alguma parte que a senhora considere passível de
        melhoria, a senhora pode relatar.
\end{enumerate}

\subsection{Tratamento dos dados}

O áudio foi transcrito e, em seguida, submetido a \textbf{edição de superfície}: supressão de
marcadores conversacionais redundantes (\textit{né}, \textit{uhum}, \textit{tá}), de
hesitações, de repetições e de trechos de conteúdo privado alheios ao objeto da pesquisa,
bem como normalização sintática. Preservaram-se integralmente o conteúdo proposicional, a
ordem dos turnos e as expressões avaliativas da participante, por constituírem o dado
primário da análise. A transcrição resultante consta da Seção~\ref{sec:transcricao}.

% =============================================================
\section{Transcrição Formalizada da Entrevista}
\label{sec:transcricao}

\noindent\textit{Convenções:} \textbf{E} designa o entrevistador (\entrevistador{}) e
\textbf{P} designa a participante (\participante{}). Intervenções de escuta ativa sem
conteúdo proposicional foram suprimidas. Colchetes indicam interpolação do pesquisador para
fins de clareza.

\begin{turnos}

\fala{E:} Dona Aline, a questão é que, depois daquela conversa que tivemos, acabamos criando
uma prévia do programa final, justamente para seguir a nossa proposta de trabalho, que é ter
uma espécie de aplicativo que atenda tanto às pessoas que precisam do serviço quanto àquelas
que podem ofertar algum tipo de serviço. Há também a possibilidade, caso a ideia seja aceita,
de firmar convênio com clínicas veterinárias e com outros tipos de profissionais, justamente
para atender essas pessoas.

\fala{E:} E também, se alguém quiser trabalhar, funcionaria basicamente como se fosse um
\textit{Uber} de mão dupla: a pessoa pode ofertar o serviço dela, como passear com o
cachorro, ficar com o cachorro e, às vezes, transportar o animal. Se a pessoa é veterinária,
ela atenderia pela plataforma do aplicativo. E aí, de acordo com a situação, isso geraria um
certo crédito para a pessoa utilizar e conseguir cuidar do bichinho.

\fala{P:} Uhum.

\fala{E:} [Demonstração no aparelho.] Ele funciona da seguinte maneira: a pessoa clica para
abrir; esta é a tela de carregamento. Há a opção de colocar usuário e senha ou de entrar pelo
Google. Quando a pessoa entra, a expectativa é vir para \enquote{Criar conta} e colocar os
dados; há também a opção de registrar pelo Google. O que conseguimos reorganizar de ontem
para hoje foi apenas a simulação do login com o Google: a pessoa entra e o programa dá um
retorno informando que está pronto para entrar.

\fala{E:} Quando a pessoa entra, ela tem a opção de escolher entre cão, gato, pássaro ou
outro. No caso de um cachorro, a pessoa coloca o nome --- o meu se chama Tobi; podemos deixar
como Tobi. Aqui, na parte de raça, é um detalhe opcional, mas vamos deixar o famoso
\enquote{sem raça definida}, o SRD.

\fala{E:} Em seguida, o sistema pede para colocar a data. Como não foi possível
\enquote{abrasileirar} o formato, por enquanto ele está no padrão americano; mas, em geral,
ele mostra o calendário, e clicando aqui fica mais fácil: escolhe-se o ano --- suponhamos que
ele tenha nascido em 2010 --- e a movimentação de mês é para a frente e para trás. Suponhamos
que tenha nascido em setembro de 2010, na mesma data de hoje. A pessoa confirma em
\enquote{OK}. Há ainda a opção de escolher uma foto na galeria ou de tirar a foto na hora.
Feito isso, o programa abre.

\fala{E:} Esta é a interface que acabamos montando. Dentro dela, há a área de cuidados do
bichinho, que mostra como está a vacinação dele, o calendário de vacinas, se ele precisa
fazer alguma coisa e se há algum agendamento. Há também a opção de ver o prontuário do
animal --- o que ele já tomou de vacina e o que ainda tem a fazer --- e a opção de editar
aqueles dados que foram cadastrados.

\fala{E:} Aqui embaixo, a pessoa alterna entre as opções: o início e a parte do animal. Esta
é a parte de triagem de serviço, ou seja, a pré-seleção do que é preciso fazer; a ideia é que
a pessoa também possa conversar com alguém por aqui. Voltando, nesta parte de rede solidária
ficam sendo mostrados os tipos de serviço que estão disponíveis; e, quando alguém pede ou
solicita algum serviço, há um mapa --- este ainda está bem genérico --- que serve para
mostrar as pessoas disponíveis.

\fala{E:} Aqui temos a configuração da pessoa em questão. Este é um nome genérico, mas
funciona da mesma forma que na parte do animal: se a pessoa quiser pesquisar na galeria ou
tirar uma foto, ela pode. E, clicando aqui ao lado, já se tem a opção de ver as informações
do programa, os termos de uso, como ele funciona, e também, como de praxe, as políticas de
privacidade, que todo programa exige.

\fala{P:} Ah, entendi.

\fala{E:} Com relação ao que a senhora viu do aplicativo, qual é a sua percepção? O que a
senhora acha?

\fala{P:} Eu acho que foi bom, é bom, e para mim é fácil. Para mim, é tipo um WhatsApp. É
basicamente isso, porque não é muito complexo.

\fala{E:} Se por acaso fosse oferecido um serviço desse tipo, seria interessante para a
senhora?

\fala{P:} Seria. Muito interessante.

\fala{E:} Então está certo. Com isso, terminamos a nossa entrevista.

\fala{P:} Mas é só isso?

\fala{E:} É só isso. É porque, dessa vez, trata-se apenas da apresentação da proposta e da
coleta de \textit{feedback}, para saber se a pessoa achou interessante. Se a senhora tiver um
ponto de vista sobre alguma parte que ache que pode melhorar, pode falar.

\fala{P:} Está muito certinho. Para mim, está bom. Para mim, é igual ao WhatsApp mesmo: está
bem prático, sem muita complicação, está bem fácil. Graças a Deus, desse jeito fica fácil de
mexer.

\subsection{Entrevista 2: Calebe (Avaliação Autônoma do Protótipo)}

\noindent\textit{Contexto:} Avaliação realizada por um jovem de 18 anos, após navegação pelo protótipo. O participante apresenta um perfil mais sistemático na análise das funções.


\fala{E:} Calebe, depois de dar uma olhada no aplicativo, o que você achou da ideia no geral?

\fala{P:} Cara, gostei muito do app, acho que pode ajudar bastante. A ideia cria uma comunidade bem legal com cada um ajudando e fazendo a sua parte, sabe? Fez sentido pra mim.

\fala{E:} Teve alguma funcionalidade específica que chamou mais a sua atenção?

\fala{P:} Sim, especialmente nessa área de triagem, onde você dá uma breve explicação do problema do pet em questão e o site já recomenda uma clínica. Achei isso bem prático e eficiente, vai direto ao ponto. E a parte dos lembretes também ajuda muito para não perder o prazo das coisas.

\fala{E:} E tem algum ponto que você achou confuso ou que a gente precisaria melhorar no app?

\fala{P:} Faltou uma aba de explicação, tipo um passo a passo para a utilização de créditos e o seu sistema. Do jeito que tá, você não entende de cara como essa moeda funciona. Tem que ter uma explicação mais sistemática de como usar isso na prática.

\end{turnos}

% =============================================================
\section{Análise dos Resultados}

\subsection{Categorização temática}

A partir das falas dos participantes, identificaram-se cinco categorias principais de análise, sintetizadas
na Tabela~\ref{tab:categorias}.

\begin{table}[h!]
\centering
\small
\begin{tabular}{@{}>{\raggedright\arraybackslash}p{3.6cm}>{\raggedright\arraybackslash}p{5.4cm}>{\raggedright\arraybackslash}p{4.4cm}@{}}
\toprule
\textbf{Categoria} & \textbf{Evidência na fala} & \textbf{Interpretação} \\
\midrule
Facilidade de uso percebida &
\enquote{Para mim é fácil}; \enquote{não é muito complexo}; \enquote{está bem fácil} &
Baixa carga cognitiva percebida; a interface não foi lida como obstáculo. \\
\addlinespace
Familiaridade e transferência de repertório &
\enquote{Para mim, é tipo um WhatsApp}; \enquote{é igual ao WhatsApp mesmo} &
A participante ancorou a compreensão do aplicativo em padrão de interação já dominado, o que
sugere aderência a convenções consolidadas de \textit{interface}. \\
\addlinespace
Intenção de uso e valor percebido &
\enquote{Seria. Muito interessante}; \enquote{está bem prático} &
Interesse declarado na adoção do serviço, associado à utilidade prática e não apenas à
estética da solução. \\
\addlinespace
Eficiência e direcionamento &
\enquote{bem prático e eficiente, vai direto ao ponto}; \enquote{ajuda muito para não perder o prazo} &
Validação direta das funcionalidades de Triagem e Lembretes como resoluções ágeis e reais para a rotina do tutor. \\
\addlinespace
Curva de aprendizado financeiro &
\enquote{Faltou uma aba de explicação [...] para a utilização de créditos}; \enquote{não entende de cara como essa moeda funciona} &
A lógica do modelo de negócios (créditos) gera atrito inicial de compreensão, evidenciando a necessidade de um \textit{onboarding} ou tutorial instrutivo. \\
\bottomrule
\end{tabular}
\caption{Categorias emergentes da análise de conteúdo das entrevistas.}
\label{tab:categorias}
\end{table}
\subsection{Discussão}

A analogia espontânea da primeira participante com um aplicativo de mensageria amplamente difundido é um achado significativo da sessão. Trata-se de indicador qualitativo de \textbf{familiaridade percebida}: não foi necessário construir um modelo mental novo para interpretar a interface. Em termos de aceitação de tecnologia, isso cobre a \textit{facilidade de uso percebida} (\enquote{é fácil}).

Complementarmente, a avaliação autônoma do segundo participante validou a \textbf{utilidade percebida} e a eficiência da proposta, com destaque para a resolutividade da \textit{Triagem Orientativa} e a conveniência dos \textit{Lembretes}. No entanto, sua navegação livre revelou um ponto de atrito crítico: a ausência de clareza sobre o funcionamento do sistema financeiro do aplicativo. Isso evidencia que, embora a interface seja familiar, o \textbf{modelo de negócios da rede solidária (créditos)} exige uma curva de aprendizado que o artefato atual não suporta de forma intuitiva sem um tutorial.

\subsection{Limitações identificadas no artefato}

Os seguintes pontos foram explicitados pelo próprio entrevistador ou identificados pelos participantes, devendo ser tratados como pendências de desenvolvimento:

\begin{itemize}[leftmargin=*]
  \item o seletor de data opera em formato norte-americano, ainda não localizado para o padrão brasileiro (\texttt{dd/mm/aaaa});
  \item o login por conta Google encontra-se \textbf{simulado}, sem integração efetiva;
  \item o mapa da rede solidária é esquemático e genérico, sem georreferenciamento real;
  \item créditos, registros e agendamentos são demonstrativos, apoiados em dados locais;
  \item \textbf{ausência de fluxos explicativos (\textit{onboarding})} ou seção de ajuda detalhando a economia e o uso dos créditos solidários.
\end{itemize}

% =============================================================
\section{Limitações Metodológicas}

O presente registro deve ser lido com as seguintes ressalvas:

\begin{enumerate}[leftmargin=*]
  \item \textbf{Amostra reduzida.} Apenas dois participantes não permitem generalização estatística; os resultados têm valor indicativo e exploratório.
  \item \textbf{Divergência na execução.} Uma interface foi manipulada pelo entrevistador (demonstração), enquanto a outra foi explorada de forma autônoma. Em ambos os casos, não houve mensuração técnica de taxa de sucesso, tempo de execução e erros de interação.
  \item \textbf{Viés de cortesia.} A existência de relação prévia entre entrevistador e participantes favorece respostas socialmente desejáveis.
  \item \textbf{Ausência de instrumento padronizado.} Não foi aplicada escala validada (como SUS ou escalas Likert de satisfação), o que impede comparação longitudinal entre versões do \textit{software}.
\end{enumerate}

% =============================================================
\section{Considerações Finais e Recomendações}

As avaliações cumpriram o objetivo de validar, junto a potenciais usuários, a inteligibilidade da proposta na versão \textit{Alpha 3i} do PetCuida. A recepção foi majoritariamente positiva quanto à facilidade de uso e ao interesse prático nas ferramentas de gestão animal. Contudo, evidenciou-se a necessidade de aprimorar a comunicação do sistema de contrapartida financeira.

Para as próximas rodadas, recomenda-se:

\begin{enumerate}[label=\textbf{R\arabic*.},leftmargin=1.6cm]
  \item ampliar a amostra, contemplando perfis distintos de tutores;
  \item padronizar a abordagem metodológica utilizando \textbf{teste de usabilidade com execução de tarefas}, registrando sucesso, hesitação e erro;
  \item aplicar instrumento padronizado ao final de cada sessão;
  \item incluir perguntas de elicitação negativa, reduzindo o efeito do viés de cortesia;
  \item priorizar, no \textit{backlog} técnico, a localização do seletor de data, a integração de autenticação e o georreferenciamento do mapa;
  \item \textbf{desenvolver uma aba explicativa ou fluxo de \textit{onboarding} focado na utilização, acúmulo e troca do sistema de créditos solidários.}
\end{enumerate}

% =============================================================
\appendix
\section{Termo de Registro e Uso dos Dados}

Os participantes foram informados sobre a finalidade acadêmica da coleta e sobre o uso restrito dos dados para fins de avaliação do artefato desenvolvido, manifestando concordância. Dados pessoais identificáveis foram suprimidos nas transcrições formalizadas.

\section{Ficha Técnica do Artefato Avaliado}

\begin{table}[h!]
\centering
\small
\begin{tabular}{@{}>{\raggedright\arraybackslash}p{4.6cm}>{\raggedright\arraybackslash}p{9.0cm}@{}}
\toprule
Denominação & PetCuida \\
Versão avaliada & \textit{Alpha 3i} \\
Plataforma & Android \\
Tecnologia & Flutter/Dart, com navegação por rotas e gerência de estado por notificadores \\
Persistência & Dados locais e demonstrativos, sem \textit{backend} \\
Escopo implementado & Autenticação, \textit{onboarding}, painel inicial, área do pet, triagem, rede solidária e perfil do tutor \\
\bottomrule
\end{tabular}
\caption{Ficha técnica resumida da versão avaliada.}
\end{table}

\section{Fontes Primárias do Registro}

\begin{itemize}[leftmargin=*]
  \item Registro em áudio da entrevista 1 (arquivo \texttt{.opus}), com duração aproximada de 7 minutos e 32 segundos;
  \item Transcrição automática com marcação temporal (arquivo \texttt{.srt}) da entrevista 1;
  \item Registros textuais de mensageria assíncrona (WhatsApp) referentes à avaliação 2;
  \item Código-fonte e documentação da versão \textit{Alpha 3i} do aplicativo.
\end{itemize}

\end{document}